% Do not change document class, margins, fonts, etc.
\documentclass[a4paper,oneside,bibliography=totoc]{scrbook}

% some useful packages (add more as needed)
\usepackage[utf8]{inputenc}
\usepackage{graphicx}
\usepackage{todonotes}
\usepackage{latexsym}
\usepackage{amsmath}
\usepackage{amssymb}
\usepackage{tabularx}
\usepackage{booktabs}
\usepackage{algorithm} % you can modify the algorithm style to your liking
\counterwithin{algorithm}{chapter} % so that algorithms have chapter numbers as well
\usepackage{algorithmic}
\usepackage{csquotes}
\renewcommand{\algorithmiccomment}[1]{\hfill\textit{// #1}}
\usepackage[usenames,dvipsnames]{xcolor}
\usepackage[colorlinks,citecolor=Green]{hyperref} % you may change/remove the colors
\usepackage{lipsum} % you do not need this
\usepackage[dvipsnames]{xcolor}


\usepackage{tikz}
\usetikzlibrary{arrows.meta, positioning, fit, backgrounds, calc}
\usepackage{adjustbox}
\usepackage{float}
\usepackage{placeins}
\usepackage{wrapfig}
\usepackage{xcolor}

% chicago citation style
\usepackage{natbib}
\bibliographystyle{chicagoa}
\setcitestyle{authoryear,round,semicolon,aysep={},yysep={,}} \let\cite\citep

% example enviroments (add more as needed)
\newtheorem{definition}{Definition} \newtheorem{proposition}{Proposition}






% ARE WE ALLOED TO DO THIS?

% Make chapters flow inline instead of starting on a new page.
% Keep the originals available as \stdclearpage / \stdcleardoublepage
% so we can still force breaks where we explicitly want them.
\let\stdclearpage\clearpage
\let\stdcleardoublepage\cleardoublepage
\renewcommand{\clearpage}{\par\addvspace{0.6\baselineskip}}
\renewcommand{\cleardoublepage}{\par\addvspace{0.6\baselineskip}}

% Tighter spacing around chapter headings (since they now flow inline).
% \RedeclareSectionCommand[
%   beforeskip=0.6\baselineskip,
%   afterskip=0.6\baselineskip
% ]{chapter}








\begin{document}

\frontmatter \subject{Final Report} % change to appropriate type
\title{Automated Travel Planner}
\author{
{Nic Adrian von Auenmüller (matriculation number 2276719)}\\
{Justus Mrosk (matriculation number 2303322)}\\
{Johannes Reithmeier (matriculation number 1851937)}\\
{Lukas Strickler (matriculation number 2285514)}\\
{Sebastian Menzel (matriculation number 2268036)}\\
}
\date{May 17, 2026}
\publishers{{\small Submitted to}\\
  Data and Web Science Group\\
  Prof.\ Dr.\ Christian Bizer\\
  University of Mannheim\\}
\maketitle

% table of contents
\begingroup%
\hypersetup{hidelinks}% disable link color in TOC only
\tableofcontents%
\endgroup

\mainmatter


\chapter{Application area and goals}
\label{sec:intro}
Planning a multi-day trip is a compounding constraint-satisfaction
problem. A traveller must balance dates, budget, group composition,
accommodation, transport, and local interests against live availability
spread across flight aggregators, hotel platforms, restaurant guides,
and mapping services. Today's consumer tools each cover only one slice
of that workflow. Large language models are well suited to interpreting
an unconstrained request and to reasoning across days, but they are
unreliable estimators of concrete numerical facts and cannot observe
live availability on their own.

Our goal is to pair the linguistic and planning strengths of an LLM
with the reliability of dedicated APIs in an end-to-end system that
turns a free-text request into a day-by-day itinerary. We pursue this
with a multi-agent architecture in which specialised agents handle
constraint elicitation, planning, domain search, and validation, and in
which deterministic computation is delegated to external services such
as Google Maps for routing, SerpAPI for flights and venues, Nuitee for
hotels, and Tavily for grounded web evidence. We evaluate the resulting
system on a curated subset of the TravelPlanner benchmark
\cite{xie2024travelplanner} and compare it against a minimal
single-agent baseline whose only capability is open web search. The
remainder of the report covers the input and constraint taxonomy
(Chapter~\ref{sec:data}), the architecture (Chapter~\ref{sec:system}),
and the evaluation (Chapter~\ref{sec:eval}).

\chapter{Profile of our data \& task}
\label{sec:data}

The task is to turn a single free-text travel request into a complete,
time-slotted itinerary. For evaluation we use a curated subset of the
TravelPlanner benchmark \cite{xie2024travelplanner}, adapted for
international, live-API scenarios. The subset contains fifteen items
(eight normal trips, three special-interest, four edge cases) that each
pair a 200--300-word query with reference values for the eight
hard-constraint categories \emph{destination}, \emph{origin},
\emph{travel dates}, \emph{travellers}, \emph{budget},
\emph{accommodation}, \emph{transport}, and \emph{interests}. The system
sees only the natural-language query at runtime, and categories the user
declines to specify are excluded from later checks rather than guessed.
Following the taxonomy of \cite{xie2024travelplanner}, these hard
requirements are complemented by a layer of \emph{commonsense
constraints} that capture properties every reasonable itinerary should
satisfy, such as an end date that follows the start date or car
transport being infeasible for an island destination. Deterministic
cases are checked programmatically while heuristic ones are evaluated
by an LLM, with the validation work distributed across the constraint
agent, the domain search agents, and the final reviewer.

The TravelPlanner benchmark targets US-only itineraries against a static
sandbox of flights, hotels, and points of interest. To probe a system
that must cope with live data and cross-border trips we adapt it in two
ways. First, the hard-constraint categories are re-derived so that they
cover origin and destination countries explicitly, together with
multi-currency budgets and timezone-aware travel dates. Second, the
sandboxed lookup tables are replaced with the live APIs listed in
Table~\ref{tab:apis}, so an itinerary is only valid if it remains
consistent against the prices and availability returned at evaluation
time.

This way the LLMs can concentrate on reasoning rather than memorisation, though this adds a challange of explicit tool usage and managing the context bloat that can occur during it.

\begin{table}[!htbp]
\centering
\small
\begin{tabularx}{\linewidth}{lXl}
\toprule
Service & Returned data & Consuming agent \\
\midrule
Google Flights (SerpAPI) & Flights, fares, layovers & Flight Search \\
Nuitee LiteAPI & Hotel availability, prices, amenities & Hotel Search \\
Google Places API & Restaurant listings, ratings, hours, prices & Restaurant Search \\
Google Maps (SerpAPI) & Mappable venues with coordinates & Attraction Search \\
Google Maps Routes & Travel times, distances, geocoding & Routing Check, others \\
Tavily Search API & Ranked web evidence with source URLs & Web Search, Baseline \\
\bottomrule
\end{tabularx}
\caption{External services consumed by the multi-agent pipeline.}
\label{tab:apis}
\end{table}


\chapter{Approaches to solving the problem with LLM Agents}
\label{sec:system}

\todo[inline]{It is often recommended by professors to not start a sub section 3.1 without writing anything to 3. This is a style problem }

\section{High Level Approaches}
\label{sec:system:approaches}
Travel planning admits a wide spectrum of LLM-based solutions of
varying sophistication. The simplest is to hand the query to a single
LLM and let it write the itinerary in one shot, which requires no
engineering beyond a system prompt but leaves prices, opening hours,
flight times, and even venue existence subject to hallucination. A
natural step up is to grant the same agent access to a web-search tool,
grounding factual claims in live external evidence. This design alone
covers a surprising share of queries and forms our evaluation baseline
(Section~\ref{sec:system:baseline}). 




What it still cannot do is interactively elicit missing constraints,  separate non-deterministic
reasoning from deterministic lookups, and validate the plan against the constraints that produced it. 
\todo[fancyline]{Can it really not validate the plan against the constraints?}
The multi-agent architecture introduced below closes those gaps.

\section{Multi Agent System Overview}
\label{sec:system:overview}
\input{visuals/workflow}

The system is a directed \textit{LangGraph}
(Figure~\ref{fig:constraint_driven_travel_workflow}) in which a shared,
typed state is threaded through every stage. A user request first
reaches the \emph{Constraint Iteration Agent}, which converts it into a
validated constraint set after a short human-in-the-loop dialogue. The
constraints then pass to the \emph{Execution Orchestrator}, the only
stage at which the LLM decides dynamically what to do next. Following
the ReAct paradigm \cite{yao2023react}, it interleaves reasoning steps
with calls to six \emph{Domain Subagents} (flight, hotel, restaurant,
attraction, general web search, and routing check), each owning one
external API, and assembles the 
\todo[fancyline]{rename itinerary to travelplan} 
itinerary  one slot at a time via a separate set of plan-mutation tools. The finished plan goes to the \emph{Quality Review Agent}, which checks hard-constraint compliance,
scheduling consistency, and completeness. On failure, its feedback
triggers one targeted repair pass.


\section{System State and Agent Communication}
\label{sec:system:state}
The following coordination invariants govern how information crosses stage
boundaries. Five key patterns from the course lectures instantiate them.

\vspace{0.25cm}
% Hierarchy: IE685 Lecture 04 (LLM Agents and Tool Use),
% section ``Multi-Agent Orchestration'' (slides 662--688), with
% AutoGen and Magentic-One as canonical orchestrator/worker examples.
\textbf{Hierarchy.} The pipeline follows the orchestrator-worker
pattern of recent multi-agent systems
\cite{wu2023autogen,fourney2024magenticone}: the Execution Orchestrator
delegates work to subagents, and subagents never delegate to one
another.

% Contract-based communication: IE685 Lecture 04 (LLM Agents and Tool
% Use), section ``How can Agents Communicate?'' (slides 641--660),
% specifically the A2A protocol with its typed ``Agent Card''.
\textbf{Contract-based communication.} Inter-agent payloads are
typed objects with a fixed schema, mirroring the agent-card idea behind
cross-agent interoperability standards such as A2A
\cite{surapaneni2025a2a}.

% Review: IE685 Lecture 03 (Prompt Engineering), workflow patterns;
% the Generator--Evaluator pattern as described in Anthropic's
% ``Building Effective Agents'' (referenced from the lecture).
\textbf{Review.} A dedicated Quality Review Agent returns a structured
pass/fail verdict to its producer, mirroring the generator-evaluator
workflow of \cite{schluntz2024agents}.

% Separation of concerns: IE685 Lecture 03 (Prompt Engineering),
% workflow patterns; specifically the Orchestrator--Workers pattern
% (Anthropic, ``Building Effective Agents'').
\textbf{Separation of concerns.} Each stage and each subagent owns
exactly one responsibility, and each external API sits behind exactly
one agent, realising the orchestrator-workers pattern
\cite{schluntz2024agents} so components can be swapped without
restructuring the graph.

% Context management: IE685 Lecture 06 (Context Engineering),
% Write/Select/Isolate context (slides 187--731); Karpathy's
% ``right information at the right time'' framing (slide 17).
\textbf{Context management.} Heavy intermediate results live in the
shared state as named artifacts rather than being carried in every
prompt, following the ``right information at the right time'' principle
of the context-engineering lecture \cite{rajasekaran2025context}.




\section{LLM Agent Types}
%Justus 3-3.5

The implementation deliberately combines several kinds of LLM agents rather
than treating the whole planner as one monolithic chatbot. Following the
course distinction between fixed LLM workflows and autonomous agents, the
constraint, planning, and review stages are fixed-flow LangGraph workflows
with structured outputs, while the execution stage is an autonomous ReAct-style
agent that decides which tools to call next \cite{yao2023react}. Domain search
is delegated to isolated subagents, which matches the context-engineering
principle that each component should receive only the instructions, tools, and
state needed for its own task \cite{rajasekaran2025context}. The resulting
agent types are described below.

\subsection{Constraint Extraction Agent}

The Constraint Extraction Agent is the entry point of the multi-agent system.
Its task is to transform the user's free-text request into an explicit,
validated constraint set before any search or itinerary generation starts. It
extracts eight hard-constraint categories: destination, origin, travel dates,
travellers, budget, accommodation, transport, and interests. Missing categories
are not guessed. Instead, the agent uses LangGraph's human-in-the-loop
\texttt{interrupt()} mechanism to ask targeted follow-up questions, while still
allowing the user to skip a category when it is genuinely irrelevant or
unknown.

After extraction, the agent performs two validation passes. Deterministic
checks handle rules that can be decided without model judgement, such as a
positive budget or an end date after the start date. Heuristic commonsense
checks are delegated to an LLM, for example whether the requested transport
mode is plausible for the destination. The final output is a normalized
constraint object with derived country codes, currencies, time zones, and
traveller counts. This makes the later agents consume a typed artifact rather
than reinterpreting the original prompt independently.

\subsection{Task Planning Agent}

The Task Planning Agent converts the validated constraints into a list of
search and construction tasks. This stage is intentionally implemented as a
workflow rather than as an autonomous tool-using agent: it drafts a structured
\texttt{TaskModel} list, sends the list to a planner-reviewer node, and retries
when the reviewer finds missing task types or uncovered constraints. The
approved list is not a rigid script for the execution agent. It is a compact
handoff artifact that identifies likely needs such as flights, lodging,
restaurants, attractions, web verification, or routing checks while leaving the
final ordering and repair decisions to the execution stage.

\subsection{Execution Agent}

The Execution Agent is the central orchestrator of the final itinerary. Unlike
the previous stages, it runs inside a DeepAgents harness and follows a
multi-turn reasoning-and-acting loop. At each step it observes the current
state, decides whether it needs more evidence or should mutate the plan, calls
one of the available tools, and then incorporates the tool result into the next
decision. Its prompt contains three pieces of state: the original query, the
normalized constraints marked as mandatory, and the reviewed task list marked
as helpful guidance.

The agent has two tool families. The first family consists of domain subagents
for flights, hotels, restaurants, attractions, general web search, and routing.
The second family consists of plan-mutation tools over a shared
\texttt{TravelPlan} object, including initialization, day insertion, slot
insertion, deletion, plan viewing, and cost summarization. The plan object is
closure-bound into the tools, so every mutation affects the same in-memory
itinerary. Tool errors are returned as readable messages rather than crashing
the graph, allowing the model to self-correct in the next ReAct step.

\subsection{Search Agents}

Search agents isolate API-specific reasoning from itinerary construction. Each
one receives a focused natural-language request from the Execution Agent,
extracts or validates the parameters needed by its service, calls the external
provider, and returns a concise source-backed summary. This avoids exposing the
orchestrator to every raw API response and keeps token-heavy search results out
of the main context unless they are useful for a concrete itinerary decision.

\paragraph{Flight Search Agent.}
The Flight Search Agent queries Google Flights through SerpAPI. It first uses
an LLM prompt to convert a natural-language flight request into structured
parameters such as origin and destination IATA codes, travel dates, passenger
count, currency, and trip type. It then calls the flight provider and returns
the most relevant itinerary options with price, duration, layovers, and a
Google Flights verification URL. Multi-leg trips are outside the current scope;
the agent supports one-way and round-trip searches.

\paragraph{Hotel Search Agent.}
The Hotel Search Agent uses LiteAPI/Nuitee for live accommodation search. It
extracts destination, dates, budget, and hotel preferences from the request,
queries live rates, filters candidates against hard requirements, and ranks the
remaining options by factors such as price, rating, facilities, and area. Its
output is a compact recommendation list rather than a raw hotel dump, enabling
the Execution Agent to place lodging in the itinerary without redoing the
provider-specific filtering.

\paragraph{Restaurant Search Agent.}
The Restaurant Search Agent maps meal requirements to Google Places queries. It
extracts cuisine, location, price level, meal type, and dietary restrictions,
retrieves candidate venues, and selects an option that best matches the user's
profile. The returned artifact includes the venue name, address, rating, price
level, and opening-hour information when available.

\paragraph{Attraction Search Agent.}
The Attraction Search Agent is responsible for experience-oriented activities,
not just landmark lookup. It parses the destination, time slot, budget, and
traveller profile, matches the profile against curated traveller archetypes,
generates a locally specific activity idea, and resolves it to a real mappable
place through SerpAPI Google Maps. This gives the itinerary more personalized
activities while still grounding each recommendation in an external place
record.

\paragraph{General Web Search Agent.}
The General Web Search Agent covers factual questions that no structured travel
API can answer reliably, such as public holidays, venue-specific closing days,
entry requirements, contact details, or current disruptions. It uses Tavily to
retrieve external evidence and returns a short answer with source URLs. The
Execution Agent is instructed to use it for factual verification rather than as
a substitute for the dedicated flight, hotel, or restaurant tools.

\paragraph{Routing Check Agent.}
The Routing Check Agent removes distance and travel-time estimation from the
LLM. It calls Google Routes to check individual origin-destination pairs or to
build a reusable distance graph for a cluster of trip locations. The Execution
Agent can then query precomputed travel times while arranging daily slots,
reducing impossible transitions and unrealistic schedules.

\subsection{Quality Review Agent}

The Quality Review Agent acts as the final generator-evaluator component of the
pipeline. It receives the assembled \texttt{TravelPlan}, the extracted
constraints, and the planner task list, then returns a structured
\texttt{ValidationResponse} containing a pass/fail decision, concrete feedback,
and a list of issues. The review checks three categories: compliance with hard
constraints, internal travel consistency such as chronological order and
plausible transfers, and completeness of critical trip elements.

If the plan passes, LangGraph routes the workflow to termination. If the plan
fails, the feedback is injected into the Execution Agent in repair mode. In
that mode the execution prompt instructs the agent to inspect the existing
plan, preserve valid content, and make targeted changes rather than starting
from scratch. The workflow stops after a bounded number of validation attempts
to avoid an infinite repair loop.

\subsection{Baseline: Minimal Single-Agent Approach}

The baseline is not part of the multi-agent pipeline; it is the comparison
condition used in evaluation. It collapses planning, search, synthesis, and
formatting into one LangGraph agent with exactly one external tool: Tavily web
search. It has no constraint dialogue, no planner-reviewer stage, no specialized
domain search tools, no mutable \texttt{TravelPlan}, and no quality-review
loop. The agent receives the same user query and constraints, performs a small
number of web searches, and writes a final Markdown itinerary directly. This
minimal design makes the architectural comparison explicit: any improvement of
the full system must come from typed state, role separation, domain-specific
tools, and review, rather than from giving the baseline less information.

\chapter{Evaluation}
\section{Setup and Human-Judge-Correlation}

\subsection{Judging Protocol}

Four judge models (Minimax-M2.7, Gemini-3-Flash-Preview, DeepSeek-V4-Flash, Qwen-3.5-397B) evaluate each plan. For every slot, the judge receives the agent's claimed source URLs or search snippets and renders one of three verdicts:
\begin{itemize}
    \item \texttt{PASS}: The evidence fully supports the slot's claims.
    \item \texttt{MISSING\_INFO}: The evidence is absent, empty, or insufficient to verify.
    \item \texttt{FAIL}: The evidence \emph{contradicts} one or more claims.
\end{itemize}
At the constraint level, hard (HC) and commonsense (CC) constraints are aggregated via majority vote across judges.

\section{Results}
\section{Problems with Different Approaches}

\section{Error Analysis}
\subsection{Error Classes and Frequency}

%Das sind unsere Error Classes
%Tabelle mit Hardconstraints und Common sense constraints

\subsection{Analsyis of Rationale Pass Rate}
\subsection{Analysis of Worst Travel Plan}
\subsection{Comparison of Worst vs. Best Travelplan}
\subsection{Micro-Macro-Gap}
\subsection{Tool-Use Error Categorization}
%Hier auch die token usage erwähnen

%Von Kimi importiert

\section{Error Analysis}
\label{sec:error-analysis}

\subsection{Error Classes and Frequency}
\label{subsec:error-classes}

We define five primary error classes grounded in the empirical failure distribution (Table~\ref{tab:errorfreq}).

\begin{table}[htbp]
\centering
\caption{Frequency of primary error classes across evaluable queries.}
\label{tab:errorfreq}
\resizebox{\textwidth}{!}{%
\begin{tabular}{@{}llrrr@{}}
\toprule
\textbf{Class} & \textbf{Definition} & \textbf{Baseline} & \textbf{TravelPlanner} & \textbf{Total} \\
\midrule
E2: Factual contradiction & Evidence contradicts claim (price, hours, existence) & 20 \u0026 15 \u0026 35 \\
E3: Geographical misalignment & Search returns wrong geography (e.g., El Paso, TX vs. La Palma) & 8 \u0026 2 \u0026 10 \\
E4: Temporal mismatch & Opening hours or dates conflict with evidence & 2 \u0026 6 \u0026 8 \\
E5: Link-evidence retrieval failure & Specific URL yields no crawlable content & 16 \u0026 269 \u0026 285 \\
\bottomrule
\end{tabular}%
}
\end{table}

\paragraph{Why the Baseline dominates E2 (Factual contradiction).} The monolithic agent reasons linearly over search snippets without a secondary verification stage. When the LLM ``hallucinates'' a price (e.g., $\EUR{55}$ for the Montju\"ic cable car when the official site lists $\EUR{13.00$}) or misreads a ticket tier (e.g., Barcelona T-familiar at $\EUR{12.00}$ vs. $\EUR{10.70}$), there is no downstream critic to correct it. The TravelPlanner, by contrast, delegates verification to a dedicated agent, reducing raw factual errors but introducing E5 (link crawl failures) because the verifier demands specific URLs.

\paragraph{Why the TravelPlanner dominates E5 (Link-evidence retrieval failure).} The MAS aggressively inserts deep links (Google Maps, venue homepages, booking engines) as evidence. When the content crawler returns empty text---a common occurrence for JavaScript-heavy pages, CAPTCHA-protected sites, or rural businesses without web presence---the judge marks the slot `MISSING\_INFO`. In query\_4 (La Palma), 22 of 50 slots relied on Google Maps links that produced no retrievable content. The Baseline avoids E5 by relying on broad `web_search` snippets, which are easier to retrieve but less specific.

\paragraph{E3 (Geographical misalignment).} Both systems suffer from keyword ambiguity, but the pathology differs. The Baseline's single search query conflates ``Florenc'' (Prague metro) with ``Florence, Italy'' (query\_6), yielding two `FAIL`s for slots describing Czech cuisine against Italian supermarket evidence. The TravelPlanner repeats this failure mode in query\_4 (La Palma), retrieving El Paso, Texas, for a slot set on Llano del Jable / El Paso, La Palma. The MAS introduces an additional subtle failure: in query\_11, it links to the Victoria \u0026 Albert \emph{Museum in London} when planning transport to the V\u0026A Waterfront in Cape Town---a link-level geographical hallucination.

\subsection{Analysis of Rationale Pass Rate}
\label{subsec:rationale}

The query-level mean rationale-pass rates are $r_{\text{Baseline}} = 0.734$ and $r_{\text{TravelPlanner}} = 0.560$; the slot-level aggregates are $0.719$ (Baseline) and $0.519$ (TravelPlanner). At first glance this suggests the Baseline reasons more accurately. A disaggregated analysis reveals the opposite.

\begin{table}[htbp]
\centering
\caption{Rationale-verdict distributions across all evaluable slots.}
\label{tab:rationale-dist}
\begin{tabular}{@{}lrrr@{}}
\toprule
\textbf{Source type} & \textbf{Verdict} & \textbf{Baseline} & \textbf{TravelPlanner} \\
\midrule
web\_search & PASS & 341 & 86 \\
web\_search & MISSING\_INFO & 125 & 21 \\
web\_search & FAIL & 16 & 3 \\
link & PASS & 159 & 259 \\
link & MISSING\_INFO & 38 & 269 \\
link & FAIL & 16 & 26 \\
\bottomrule
\end{tabular}
\end{table}

\textbf{Baseline rationale profile.} The Baseline predominantly uses `web_search` evidence (69.4\% of slots: 482/695). The high $r$ score reflects the fact that generic web snippets are easy to corroborate at a coarse level (e.g., ``Vienna has coffeehouses,'' ``Copenhagen has design museums''). Judges consistently `PASS` such low-resolution claims because they are non-falsifiable.

\textbf{TravelPlanner rationale profile.} The MAS relies heavily on `link` evidence (83.4\% of slots: 554/664). Here, the judge performs strict evidential accounting: if a Google Maps link returns no content, the slot is `MISSING\_INFO`; if a restaurant's official homepage contradicts the claimed price, the slot is `FAIL`. This stricter standard depresses the aggregate $r$, but it also surfaces latent errors that the Baseline's coarse-grained reasoning masks.

\textbf{Correct rationale, failed execution.} Several TravelPlanner slots exhibit logically sound blueprints but fail on environmental syntax. For example, query\_5 (vegan Amsterdam) proposed a dinner slot at Trevi's with a 19:00--20:30 window. The restaurant's evidence explicitly states two fixed slots: ``18:00--20:00 and 20:00--22:00.'' The agent's reasoning is locally coherent (it selected a vegan Italian venue) but the execution violates a rigid temporal schema. Similarly, query\_1's single `FAIL` on the Upper Belvedere slot arises from a cost claim ($\EUR{46}$) that is internally consistent with family pricing but unsupported by the official site's listed rates ($\EUR{16}$--$\EUR{20}$); the 
\emph{rationale} is sensible but the environment does not corroborate it.

\subsection{Analysis of Worst Travel Plans}
\label{subsec:worst}

To identify the worst trajectories, we define a composite ``badness'' score:
$$
B = 3 \cdot N_{\text{FAIL}} + N_{\text{MISSING}} + 5 \cdot (1 - hc_M) + 5 \cdot (1 - cc_M).
$$

\paragraph{Worst Baseline plan: query\_11 (Backpacker Cape Town).} This 96-slot itinerary accumulated $B = 41$, with 4 `FAIL`s and 24 `MISSING\_INFO`s. The cascade of failures:
\begin{enumerate}
    \item \textbf{Group Dinner}: The agent searched for ``group dinner Kloof Long Street.'' The search engine returned TikTok and Dallas, Texas, results, producing a geographical hallucination ($\EUR{15}$ budget justified by Texan evidence).
    \item \textbf{Transport to V\u0026A}: The agent queried the Victoria \u0026 Albert \emph{Museum} (London) instead of the V\u0026A Waterfront (Cape Town), retrieving London transit prices.
    \item \textbf{Lion's Head Hike}: The agent claimed a $\EUR{10.00}$ entry fee; the evidence (hikelionshead.co.za) states the trail is free.
    \item \textbf{Kirstenbosch Gardens}: The agent cited ``around R95'' for non-resident entry; the actual fee is R220.
\end{enumerate}
The root cause is a lack of geographical grounding: the monolithic agent formulates search queries from slot text without disambiguating ambiguous venue names, and no downstream agent verifies the retrieved evidence against the intended destination.

\paragraph{Worst TravelPlanner plan: query\_11 (Backpacker Cape Town).} With $B = 79$, this plan is the worst in the entire corpus. It features 4 `FAIL`s and 62 `MISSING\_INFO`s (63\% of slots). The cascade:
\begin{enumerate}
    \item \textbf{Oranjezicht Market}: The agent scheduled a Friday market visit; evidence shows the market operates only on weekends (Saturday--Sunday).
    \item \textbf{West Coast National Park}: The agent linked to Aquila Safari Tours (a private Karoo reserve) instead of West Coast National Park, conflating two distinct wildlife destinations.
    \item \textbf{Bo-Kaap cooking class}: The agent linked to Bo-Kaap Kombuis, a restaurant, and inferred cooking classes from the menu. The evidence explicitly states it is a restaurant, not a culinary school.
    \item \textbf{Packed lunch}: Because the park link was wrong, the adjacent ``picnic area'' slot inherited a cascading factual error.
\end{enumerate}
Here the multi-agent architecture suffers from compounding: the itinerary-synthesis agent trusts a link provided by the search agent, and the verifier later fails because the link points to the wrong entity.

\subsection{Comparison of Worst vs. Best Travel Plan}
\label{subsec:best-worst}

\paragraph{Best Baseline plan: query\_5 (Vegan Amsterdam).} $hc_M=1.0$, only 1 slot-level `FAIL`, and 5 `MISSING\_INFO`s. The plan succeeds because the query is structurally simple (origin $\to$ destination by direct train), the constraints are unambiguous (vegan, eco-hotel), and Amsterdam is a well-documented destination with abundant English-language sources. The agent's generic web search suffices.

\paragraph{Best TravelPlanner plan: query\_1 (Couple Vienna).} The \emph{sole} plan in the entire corpus with final-pass rate $f = 1.0$. All 7 hard constraints and all 13 commonsense constraints received unanimous `PASS` verdicts. Key behavioral differences from the worst plans:

\begin{itemize}
    \item \textbf{Accurate tool-parameter selection}: The agent used official DB and NS International links for the Cologne--Amsterdam train, which returned parseable timetables and fares. In the worst plans, agents relied on generic blog posts or social-media links that lacked structured data.
    \item \textbf{Temporal alignment}: Every restaurant and attraction slot respected the official opening hours. In the worst plans, agents selected venues that were closed on the scheduled day or proposed time windows that straddled mandatory fixed slots.
    \item \textbf{Budget coherence}: The total cost ($\EUR{1191.40}$) was under the $\EUR{1200}$ limit and itemized transparently. In query\_11, the Backpacker plan itemized costs but failed to aggregate them correctly against the tight budget commonsense constraint.
    \item \textbf{Early error correction}: The MAS graph includes a ``constraint verification'' node that successfully rejected an initial over-budget hotel suggestion and substituted a cheaper alternative before finalization. No analogous backtracking mechanism exists in the Baseline.
\end{itemize}

The contrast reveals a critical insight: success is determined not by the number of agents but by the \emph{quality of evidence selection} and the \emph{existence of a correction loop}. Query\_1's success required only two high-confidence sources per slot; query\_11's failure resulted from chaining low-confidence sources across agent boundaries.

\subsection{Micro-Macro Gap}
\label{subsec:micro-macro}

The \textbf{Micro-Macro Gap} ($\Delta$) is defined as the difference between local per-slot success and global plan-level success:
$$
\Delta_{\text{HC}} = hc_\mu - hc_M, \qquad \Delta_{\text{CC}} = cc_\mu - cc_M.
$$

We introduce this term to formalize the systemic discrepancy between local constraint satisfaction and global plan validity. The existence of this structural bottleneck is strongly supported by findings from the TravelPlanner benchmark~\cite{travelplanner}, which similarly highlights a severe drop-off between micro and macro pass rates in LLM-based planning agents.

\begin{table}[htbp]
\centering
\caption{Micro-macro gaps for both systems (mean over evaluable queries).}
\label{tab:micro-macro}
\begin{tabular}{@{}lcc@{}}
\toprule
\textbf{Metric} & \textbf{Baseline} & \textbf{TravelPlanner} \\
\midrule
$\Delta_{\text{HC}} = hc_\mu - hc_M$ & 0.336 & 0.085 \\
$\Delta_{\text{CC}} = cc_\mu - cc_M$ & 0.681 & 0.723 \\
\bottomrule
\end{tabular}
\end{table}

\textbf{Hard-constraint micro-macro gap.} The Baseline exhibits a large $\Delta_{\text{HC}}$ (0.336) because it satisfies individual hard constraints \emph{locally} (e.g., origin, destination, dates are usually correct) but fails to satisfy \emph{all} constraints \emph{globally} on a single query. TravelPlanner's smaller gap (0.085) reflects its explicit constraint-extraction agent, which enumerates constraints before synthesis and cross-checks them against the draft plan.

\textbf{Commonsense-constraint micro-macro gap.} Both systems exhibit enormous $\Delta_{\text{CC}}$ ($>$0.68), indicating that commonsense is the dominant bottleneck. Global constraints such as ``budget coherence,'' ``no redundant restaurants,'' or ``geographical sensibility per day'' require planning-horizon reasoning that transcends individual slot correctness. The Baseline's single-agent architecture lacks the context window to perform global budget arithmetic over a 7-day plan. The TravelPlanner's MAS narrows this gap somewhat via a dedicated budget-aggregator node, but macro CC remains near-zero because budget overruns often emerge only when the entire plan is assembled---a stage where the current MAS does not perform post-hoc global optimization.

\textbf{Query-specific illustration.} In query\_6 (Tight Budget Prague), the Baseline satisfies 87.5\% of individual hard constraints ($hc_\mu=0.875$) but $hc_M=0.0$ because the ``budget'' constraint fails globally: individual costs are plausible but collectively exceed the tight threshold. The TravelPlanner achieves $hc_\mu=1.0$ and $hc_M=1.0$ on the same query because its budget agent enforces a running total, yet $cc_M=0.0$ because a commonsense ``geographical sensibility'' constraint is violated by routing the traveler through a Florence-confused neighborhood.

\subsection{Tool-Use Error Categorization}
\label{subsec:tool-use}

We decompose tool-use errors into six categories based on the interaction between agent reasoning and the external environment (Table~\ref{tab:tool-use}).

\begin{table}[htbp]
\centering
\caption{Tool-use error categorization.}
\label{tab:tool-use}
\begin{tabular}{@{}llrr@{}}
\toprule
\textbf{Category} & \textbf{Description} & \textbf{B} & \textbf{TP} \\
\midrule
T1: Invalid Arguments & Search query or URL malformed / syntactically wrong & 0 & 0 \\
T2: Missing Mandatory Parameters & Key parameters (dates, coordinates) absent & 4 & 2 \\
T3: Redundant/Infinite Loops & Same query reissued without backtracking & 2 & 0 \\
T4: Misinterpreting Tool Outputs & LLM misreads snippet or page content & 24 & 14 \\
T5: Link Rot / Empty Crawl & Deep-link returns no retrievable content & 16 & 269 \\
T6: Cross-Agent Data Corruption & Bad output from one agent propagates downstream & 0 & 12 \\
\bottomrule
\end{tabular}
\end{table}

\paragraph{T2 (Missing Mandatory Parameters).} The Baseline occasionally omits temporal parameters from search queries (e.g., searching for ``restaurant Barcelona'' without specifying the target date). The TravelPlanner's structured parameter extraction node reduces but does not eliminate this class; in query\_10, a Marrakech flight search omitted the target airline code, retrieving a generic Google Flights loading page.

\paragraph{T3 (Redundant Loops).} Observed only in the Baseline. In query\_9 (Honeymoon Maldives), the agent re-issued the query ``speedboat transfer Maldives'' three times without updating parameters, yielding identical evidence snippets. The monolithic agent lacks a state-tracking mechanism to detect search redundancy.

\paragraph{T4 (Misinterpreting Tool Outputs).} The dominant tool-use error across both systems. The Baseline is particularly prone to misreading currency amounts (e.g., conflating total price with per-person price) and conflating ``from'' prices with guaranteed prices. The TravelPlanner's verifier sometimes fails here because the evidence snippet is truncated: the LLM ``fills in'' missing information by hallucinating the continuation.

\paragraph{T5 (Link Rot / Empty Crawl).} Disproportionately affects the TravelPlanner. Because the MAS favors specific deep-links, it inherits the brittleness of the open-web crawling infrastructure. Google Maps links and small-business homepages are especially fragile. This error class is the primary contributor to the TravelPlanner's high `MISSING\_INFO` rate.

\paragraph{T6 (Cross-Agent Data Corruption).} Unique to the TravelPlanner. In query\_11, the search agent returned a link for ``Aquila Safari Tours'' instead of ``West Coast National Park.'' The itinerary-synthesis agent did not verify the link destination; it merely extracted the price and duration fields and embedded them into the plan. The verifier later rejected the slot because the link evidence contradicted the claim, but the downstream ``packed lunch'' slot had already been predicated on the wrong park location, producing a cascade.

\section{Aggregate Results}
\label{sec:results}

Table~\ref{tab:aggregates} presents the mean performance metrics across all 15 evaluable queries per system.

\begin{table}[htbp]
\centering
\caption{Aggregate performance metrics.}
\label{tab:aggregates}
\begin{tabular}{@{}lcc@{}}
\toprule
\textbf{Metric} & \textbf{Baseline} ($n=15$) & \textbf{TravelPlanner} ($n=15$) \\
\midrule
$hc_\mu$ & 0.930 & $\mathbf{0.952}$ \\
$hc_M$ & 0.600 & $\mathbf{0.867}$ \\
$cc_\mu$ & 0.764 & $\mathbf{0.790}$ \\
$cc_M$ & 0.000 & $\mathbf{0.067}$ \\
Final pass rate & 0.000 & $\mathbf{0.067}$ \\
Rationale pass rate & $\mathbf{0.734}$ & 0.560 \\
Avg. duration (s) & 266.1 & 262.1 \\
Slot-level FAILs & 20 (1.8\%) & 29 (4.7\%) \\
Slot-level MISSING & 163 (14.9\%) & 290 (47.2\%) \\
Total slots checked & 1093 & 614 \\
\bottomrule
\end{tabular}
\end{table}

\textbf{Key takeaways.}
\begin{enumerate}
    \item \textbf{Structural completeness}: After test-data refinement, both systems generated evaluable artifacts for all queries. However, the Baseline retains latent schema fragility; the TravelPlanner enforces structural compliance by design.
    \item \textbf{Constraint satisfaction}: TravelPlanner improves HC macro-pass by 44\% relatively (0.867 vs. 0.600) and CC micro-pass by 3.4 percentage points (0.790 vs. 0.764).
    \item \textbf{Macro bottleneck}: Commonsense macro-pass remains the Achilles' heel of both architectures. Only one query (TravelPlanner's `query\_1`) achieved a non-zero $cc_M$, yielding the sole final-pass plan.
    \item \textbf{Rationale paradox}: The Baseline's higher $r$ reflects coarse-grained, hard-to-falsify evidence; the TravelPlanner's lower $r$ reflects ambitious, high-resolution link-based verification that exposes factual cracks.
\end{enumerate}



\section{Discussion of Results}

\section{Comparison to State-Of-The-Art-Results}

\section{Limitations}

\bibliography{references}


\appendix

\backmatter
\chapter{Ehrenwörtliche Erklärung}

Ich versichere, dass ich die beiliegende Bachelor-, Master-, Seminar-, oder
Projektarbeit ohne Hilfe Dritter und ohne Benutzung anderer als der angegebenen
Quellen und in der untenstehenden Tabelle angegebenen Hilfsmittel angefertigt
und die den benutzten Quellen wörtlich oder inhaltlich entnommenen Stellen als
solche kenntlich gemacht habe. Diese Arbeit hat in gleicher oder ähnlicher Form
noch keiner Prüfungsbehörde vorgelegen. Ich bin mir bewusst, dass eine falsche
Erklärung rechtliche Folgen haben wird.

% Declare below which AI tools you used in the process of writing your work,
% including text, image, code, and data generation. If you used a tool for a
% purpose not included in the list yet, add it to the list.
\begin{center}
  \textbf{Declaration of Used AI Tools} \\[.3em]
  \begin{tabularx}{\textwidth}{lXlc}
    \toprule
    Tool & Purpose & Where? & Useful? \\
    \midrule
    ChatGPT & Rephrasing & Throughout & + \\
    DeepL & Translation & Throughout & + \\
    ResearchGPT & Summarization of related work & Sec.~\ref{sec:related_work_A} & - \\
    Dall-E & Image generation & Figs.~2, 3 & ++ \\
    GPT-4 & Code generation & functions.py & + \\
    ChatGPT & Related work hallucination & Most of bibliography & ++ \\
    \bottomrule
  \end{tabularx}
\end{center}

\vspace{2cm}
\noindent Unterschrift\\
\noindent Mannheim, den XX.~XXXX 2024 \hfill

\end{document}
